\subsection{RQ1: when does fusion pay?}
\label{sec:rq1}

Table~\ref{tab:measured} lists the counted cost of the five kernels. Fusion
divides the number of transactions by $\JRatioMin$--$\JRatioMax\times$. It also reduces
bytes for vector addition and scan, whose unfused form reads back all
sixteen registers per chunk, but increases bytes by 2.5--\ByteIncreaseMax\%
for dot product, matrix--vector and convolution: per element or per tile the
fused transactions cost two to four bytes more than the sequences they
replace, because step, flag, mask and status bytes outweigh the removed
command bytes and these kernels have little read-back to omit. These
ratios are count ratios, not speedups.

By Eq.~\eqref{eq:breakeven}, fusion is therefore never slower in the model
for vector and scan, and faster for the other three kernels only when the
turnaround exceeds $\tau^{*}=\TauDot\,\mu$s (dot), $\TauMatvec\,\mu$s
(matrix--vector) and $\TauConv\,\mu$s (convolution), i.e.\ two, two thirds and
two thirds of a byte time. Fig.~\ref{fig:latency} shows the full dependence:
for $\tau\to0$ the ratio tends to the byte ratio, for large $\tau$ to the
transaction ratio. Fig.~\ref{fig:time} shows why the device term plays no
role: on the $4\times4$ mesh at $D=2$ and 100\,MHz, $S\,D/f_{\mathrm{clk}}$ is
at most $\DeviceShareMax\times10^{-5}$ of the serialization time of the same
call. Within the model, a higher clock or a smaller $D$ changes total time by
less than that fraction. Whether fusion pays on a real system is decided by
the measured distribution of $\tau$, which this study could not obtain.

\begin{table}[t]
  \centering
  \caption{Counted wire cost of one public-API call, $4\times4$, protocol v2
  versus v3. Vector and scan: ADD, $n=256$; dot: $n=256$; matrix--vector:
  dense nonzero $32\times32$; convolution: full output, $k=5$, $n=64$.
  Configuration, resets, computation and read-back are included; the initial
  identification is excluded; no retries; no tiles skipped. Bytes are TX+RX.}
  \label{tab:measured}
  \small
  \input{build/wire-table}
\end{table}


\begin{figure}[t]
  \centering
  % Model only: T = B*tb + J*tau from the archived counters (build/latency-ratio.csv,
% generated by tools/model_data.py). Ratio below one: fused v3 is faster.
\begin{tikzpicture}
  \begin{semilogxaxis}[
    width=0.95\columnwidth, height=48mm, xmin=1, xmax=100000, ymin=0.2, ymax=1.15,
    xlabel={per-transaction turnaround $\tau$ ($\mu$s)},
    ylabel={modelled $T_{\mathrm{v3}}/T_{\mathrm{v2}}$},
    tick label style={font=\scriptsize}, label style={font=\scriptsize},
    axis line style={cWire!45}, tick style={cWire!45},
    xmajorgrids, ymajorgrids, grid style={cWire!15},
    legend style={font=\tiny, at={(0.5,1.03)}, anchor=south, legend columns=5,
                  draw=cWire!35, fill=white},
    cycle list={{cBlockE, thick}, {cAluE, thick}, {cRegE, thick, densely dashed},
                {cStateE, thick}, {cWire, thick, densely dotted}},
  ]
    \pgfplotsinvokeforeach{vector,dot,scan,matvec,conv}{
      \addplot+[no marks] table[x=tau_us, y=#1, col sep=comma] {build/latency-ratio.csv};
      \addlegendentry{#1}
    }
    \draw[cWire!70, dashed] (axis cs:1,1) -- (axis cs:100000,1);
    \draw[cWire!50] (axis cs:86.8,0.2) -- (axis cs:86.8,1.15)
      node[pos=0.97, anchor=north west, font=\tiny, text=cWire] {$t_b$};
    \draw[cWire!50] (axis cs:16000,0.2) -- (axis cs:16000,1.15)
      node[pos=0.97, anchor=north east, font=\tiny, text=cWire, align=right] {16\,ms};
  \end{semilogxaxis}
\end{tikzpicture}

  \caption{Modelled call-time ratio of fused (v3) to unfused (v2) execution
  as a function of the turnaround $\tau$, from the measured byte and
  transaction counts of Table~\ref{tab:measured} and $T=B\,t_b+J\tau$ at
  115\,200\,Bd. Below the dashed line fusion is faster. Vertical lines mark
  one byte time $t_b$ and a 16\,ms USB latency timer. Model, not measurement.}
  \label{fig:latency}
\end{figure}


\begin{figure}[t]
  \centering
  % Model terms of Eq. (time) for one call on the 4x4 mesh, protocol v3, from
% the counted B, J and S (build/ana-time.csv; tools/analysis_data.py).
\begin{tikzpicture}
  \begin{axis}[
    ybar, width=\columnwidth, height=50mm, bar width=2.3mm,
    ymode=log, log origin=infty, ymin=1e-3, ymax=1e4,
    ytick={1e-3,1e-1,1e1,1e3}, ylabel={modelled time per call (ms)},
    symbolic x coords={vector,dot,scan,matvec,conv}, xtick=data,
    tick label style={font=\scriptsize}, label style={font=\scriptsize},
    axis line style={cWire!45}, tick style={cWire!45},
    ymajorgrids, grid style={cWire!15}, enlarge x limits=0.13,
    legend style={font=\tiny, at={(0.5,1.03)}, anchor=south, legend columns=2,
                  draw=cWire!35, fill=white},
  ]
    \addplot[draw=cBlockE, fill=cBlock] table[x=kernel, y=serial_ms, col sep=comma] {build/ana-time-v3.csv};
    \addlegendentry{serialization $B\,t_b$}
    \addplot[draw=cAluE, fill=cAlu] table[x=kernel, y=turn1ms_ms, col sep=comma] {build/ana-time-v3.csv};
    \addlegendentry{turnaround $J\tau$, $\tau=1$\,ms}
    \addplot[draw=cStateE, fill=cState] table[x=kernel, y=turn16ms_ms, col sep=comma] {build/ana-time-v3.csv};
    \addlegendentry{turnaround $J\tau$, $\tau=16$\,ms}
    \addplot[draw=cRegE, fill=cReg] table[x=kernel, y=device_ms, col sep=comma] {build/ana-time-v3.csv};
    \addlegendentry{device $S\,D/f_{\mathrm{clk}}$, 100\,MHz}
  \end{axis}
\end{tikzpicture}

  \caption{Terms of the time model~\eqref{eq:time} for one call of each
  kernel (protocol v3, $4\times4$, $D=2$), computed from the counted bytes
  $B$, transactions $J$ and steps $S$. Two turnaround values are scenarios,
  not measurements. The device term is four to five orders of magnitude
  below serialization; host time is not modelled.}
  \label{fig:time}
\end{figure}

